% from ega0-20-fr.tex
\oldpage[0\textsubscript{IV-1}]{144}
\begin{corollary}[20.6.20]
\label{0.20.6.20-fr}
\begin{enumerate}
  \item[\rm(i)] Si $K$ est une extension algébrique séparable de $k$, on a
  $\Omega^1_{K/k}=0$.
  \item[\rm(ii)] Soient $k$ un corps de caractéristique 0, $K$ une extension
  de $k$. Pour que $\Omega^1_{K/k}=0$, il faut et il suffit que $K$ soit une
  extension algébrique de $k$. En particulier, pour qu'un corps $K$ de
  caractéristique 0 soit tel que $\Omega^1_K=0$, il faut et il suffit que
  $K$ soit une extension algébrique de $\mathbb{Q}$ (cf.\ (21.4.4) et
  (21.7.5)).
\end{enumerate}
\end{corollary}

\begin{proof}
(i) Pour tout $x\in K$, soit $f$ le polynôme minimal de $x$ sur $k$ ; comme
$f'(x)\ne0$ et $d_{K/k}(f(x))=f'(x)d_{K/k}(x)=0$, on a
$d_{K/k}(x)=0$ et notre assertion résulte de \sref{0.20.4.7-fr}.

(ii) Il existe une extension pure $L$ de $k$ telle que
$k\subset L\subset K$ et que $K$ soit une extension algébrique de $L$.
Comme $K$ est séparable sur $L$, il résulte de \sref{0.20.6.19-fr}, (i) que
la suite \sref{0.20.5.7.2-fr}
\[
  0\to\Omega^1_{L/k}\otimes_L K\to\Omega^1_{K/k}
  \to\Omega^1_{K/L}\to0
\]
est exacte et de (i) que $\Omega^1_{K/L}=0$. La relation
$\Omega^1_{K/k}=0$ est donc équivalente à $\Omega^1_{L/k}=0$, et comme
$L$ est une extension pure de $k$, il résulte de \sref{0.20.5.10-fr} que
la relation $\Omega^1_{L/k}=0$ équivaut à $L=k$.
\end{proof}

\begin{remarks}[20.6.21]
\label{0.20.6.21-fr}
\begin{enumerate}
  \item[\rm(i)] Comme $\Upsilon_{B/A/\Lambda}$ est le noyau de
  \[
    u_{B/A/\Lambda}:\Omega^1_{A/\Lambda}\otimes_A B
    \to\Omega^1_{B/\Lambda}
  \]
  et $\Upsilon^C_{B/A/\Lambda}$ le noyau de
  $u_{B/A/\Lambda}\otimes1_C$, on a un homomorphisme canonique
\begin{equation}
\tag{20.6.21.1}
\label{0.20.6.21.1-fr}
  \Upsilon_{B/A/\Lambda}\otimes_B C\to\Upsilon^C_{B/A/\Lambda}.
\end{equation}
Cet homomorphisme est bijectif lorsque la suite
\begin{equation}
\tag{20.6.21.2}
\label{0.20.6.21.2-fr}
  0\to\Upsilon_{B/A/\Lambda}\otimes_B C\to
  \Omega^1_{A/\Lambda}\otimes_A C\to
  \Omega^1_{B/\Lambda}\otimes_B C\to
  \Omega^1_{B/A}\otimes_B C\to0
\end{equation}
est exacte, ce qui se produit dans les cas suivants :
\begin{enumerate}
  \item[$1^\circ$] $C$ est un $B$-module \emph{plat}.
  \item[$2^\circ$] Les $B$-modules $\Omega^1_{B/\Lambda}$ et
  $\Omega^1_{B/A}$ sont \emph{plats} : en effet, il en est alors de même de
  $\operatorname{Ker}(\Omega^1_{B/\Lambda}\to\Omega^1_{B/A})$
  ($\mathbf{0_I}$, 6.1.2), et la suite \sref{0.20.6.21.2-fr} est alors
  exacte en vertu de ($\mathbf{0_I}$, 6.1.2).
\end{enumerate}
  \item[\rm(ii)] Considérons un diagramme commutatif d'homomorphismes
  d'anneaux
\[
\xymatrix{
  \Lambda' \ar[r] & A' \ar[r] & B' \ar[r] & C' \\
  \Lambda \ar[u] \ar[r] & A \ar[u] \ar[r] & B \ar[u] \ar[r] & C \ar[u]
}
\]
Alors les définitions de \sref{0.20.6.16-fr} montrent que l'on a un
diagramme commutatif de complexes (où les flèches verticales proviennent de
\sref{0.20.6.6.4-fr})
\[
\xymatrix{
  0 \ar[r] & \mathrm{K}^C_\bullet(B/A/\Lambda) \ar[r] \ar[d] &
  \mathrm{K}'_\bullet(C/A/\Lambda) \ar[r] \ar[d] &
  \mathrm{K}_\bullet(C/B/\Lambda) \ar[r] \ar[d] & 0 \\
  0 \ar[r] & \mathrm{K}^{C'}_\bullet(B'/A'/\Lambda') \ar[r] &
  \mathrm{K}'_\bullet(C'/A'/\Lambda') \ar[r] &
  \mathrm{K}_\bullet(C'/B'/\Lambda') \ar[r] & 0
}
\]
